\clearpage
\begin{proofbox}
\textbf{\textcolor{CProof}{方法一：最快证明（考场首选｜设线 + 韦达，约三步）}}\par\smallskip
\textbf{思路：}直线过焦点，先设 $x=my+\frac p2$；\textbf{只看二次方程的常数项}，再由抛物线方程得到横积。
\[
 \begin{aligned}
 x=my+\frac p2,\quad y^2=2px
 &\Longrightarrow y^2-2pmy-p^2=0,\\
 &\overset{\text{韦达}}{\Longrightarrow}
 \boxed{y_1y_2=-p^2},\\[2pt]
 x_i=\frac{y_i^2}{2p}
 &\Longrightarrow
 \boxed{x_1x_2=\frac{(y_1y_2)^2}{4p^2}=\frac{p^2}{4}}.
 \end{aligned}
\]
\textbf{一句话抓手：}过焦点 $\Rightarrow$ 方程常数项为 $-p^2$
$\Rightarrow$ 纵积固定 $\Rightarrow$ 横积固定。
\end{proofbox}
\clearpage
\begin{proofbox}
\textbf{\textcolor{CProof}{方法二：最低门槛，相似三角形加因式分解}}

\textbf{先沿用前两幅图。}非竖直且正斜率时，把右上端记为 $A$、左下端记为 $B$。令 $u=y_A>0$、$v=-y_B>0$，焦点横坐标为 $p/2$。
\[
 x_A=\frac{u^2}{2p},\qquad x_B=\frac{v^2}{2p},\qquad
 h=x_A-\frac p2,\qquad k=\frac p2-x_B.
\]
由于焦点在两端点之间，$h,k>0$。\textbf{第1步，只用相似：}
\[
 \frac hu=\frac kv.
\]
\textbf{第2步，用抛物线方程替换水平边：}
\[
 \frac{u^2-p^2}{2pu}=\frac{p^2-v^2}{2pv}
 \quad\Longrightarrow\quad v(u^2-p^2)=u(p^2-v^2).
\]
\textbf{第3步，移项并提取共同因子：}
\[
 u^2v+uv^2=p^2(u+v)
 \quad\Longrightarrow\quad (u+v)(uv-p^2)=0.
\]
$u+v>0$，故 $uv=p^2$。注意 $B$ 的纵坐标是 $-v$，所以
\[
 \boxed{y_Ay_B=-uv=-p^2},\qquad
 \boxed{x_Ax_B=\frac{u^2v^2}{4p^2}=\frac{p^2}{4}}.
\]
\textbf{补齐情况。}负斜率时先上下翻折，边长与横积不变；两个纵坐标同时变号，纵积也不变。竖直通径单独代入 $x=p/2$，得 $y=\pm p$，两个积同样成立。

\medskip
\textbf{也可直接用带符号的坐标，不作图形分类。}设 $A(a^2/(2p),a)$、$B(b^2/(2p),b)$，其中 $a,b\ne0$ 且 $a\ne b$。三点共线给
\[
 \frac{a^2-p^2}{2pa}=\frac{b^2-p^2}{2pb}
 \quad\Longrightarrow\quad (a-b)(ab+p^2)=0.
\]
约去 $a-b\ne0$，仍得 $ab=-p^2$。这就是前面相似比的带符号写法；不需要韦达定理。
\end{proofbox}
\clearpage
\begin{proofbox}
\textbf{\textcolor{CProof}{方法三：常规证明（点斜式 + 联立 + 韦达｜补齐竖直特例）}}\par\smallskip
\textbf{情况 1：}直线斜率存在。设 $y=k(x-\frac p2)$。
若 $k=0$，直线就是 $y=0$，它只交抛物线于顶点，不符合两个不同交点的条件，故此处 $k\ne0$。
联立得到关于 $x$ 的方程：
\[
 k^2\left(x-\frac p2\right)^2=2px
 \Longrightarrow
 k^2x^2-(pk^2+2p)x+\frac{p^2k^2}{4}=0.
\]
韦达定理直接给出
\[
 x_1+x_2=p+\frac{2p}{k^2},\qquad
 \boxed{x_1x_2=\frac{p^2}{4}}.
\]
再由 $y_i=k(x_i-\frac p2)$，
\[
 \begin{aligned}
 y_1y_2&=k^2\left[x_1x_2-\frac p2(x_1+x_2)+\frac{p^2}{4}\right]\\
 &=k^2\left[\frac{p^2}{4}-\frac p2\left(p+\frac{2p}{k^2}\right)+\frac{p^2}{4}\right]
 =\boxed{-p^2}.
 \end{aligned}
\]
\textbf{情况 2：}直线斜率不存在。这就是通径 $x=\frac p2$。
代入 $y^2=2px$ 得 $y^2=p^2$，故
\[
 A\!\left(\frac p2,p\right),\quad
 B\!\left(\frac p2,-p\right),\qquad
 x_1x_2=\frac{p^2}{4},\quad y_1y_2=-p^2.
\]
两种情况合起来，证明完成。
\end{proofbox}
\clearpage
\begin{proofbox}
\textbf{\textcolor{CProof}{推广证明：切线垂直，且交点在准线上}}\par\smallskip
对 $y^2=2px$，在 $P(x_i,y_i)$ 处的切线为
\[
 y_iy=p(x+x_i),\qquad k_i=\frac p{y_i}.
\]
由于 $y_1y_2=-p^2$，两条切线的斜率积为
\[
 \boxed{k_Ak_B=\frac{p^2}{y_1y_2}=-1}.
\]
再设两切线交于 $T(X,Y)$。相减两切线方程：
\[
 (y_1-y_2)Y=p(x_1-x_2)
 =\frac{y_1^2-y_2^2}{2}.
\]
约去 $y_1-y_2\ne0$，得 $Y=\frac{y_1+y_2}{2}$。
代回 $y_1Y=p(X+x_1)$，用 $px_1=\frac{y_1^2}{2}$，得到
\[
 pX=\frac{y_1(y_1+y_2)}{2}-\frac{y_1^2}{2}
 =\frac{y_1y_2}{2}=-\frac{p^2}{2}.
\]
因此
\[
 \boxed{T\!\left(-\frac p2,\frac{y_1+y_2}{2}\right)},
\]
横坐标等于准线 $x=-\frac p2$；证毕。

\medskip
\textbf{\textcolor{CProof}{证明方法怎么选？}}\par\smallskip
\textbf{求快}选方法一；\textbf{想看清为什么成立、不想碰韦达}选方法二；
\textbf{题目已用 $y=k(x-\frac p2)$ 设线}选方法三，并记得补竖直通径。
\end{proofbox}
